% SOURCE_AUTHORITY: sga5_fr_workpass.tex, lines 6794--6859 (LF-normalized unit).
% SOURCE_SHA256: 791F4EFFC5E02832D5D77ED03518C8156D6F07E4C8238B03545DB93D883FBB28
\subsubsection*{2.3. Lembretes sobre o cálculo de frações e aplicações}
Sejam $\cE$ uma categoria e $S$ um subconjunto de $\Fl(\cE)$. Chama-se categoria de frações de $\cE$ em relação a $S$ um par $(p,\cE(S^{-1}))$ formado por uma categoria $\cE(S^{-1})$ e um funtor $p:\cE\to\cE(S^{-1})$ que verifica a seguinte propriedade universal: para que um funtor $u$ de $\cE$ em uma categoria $\cF$ se fatore através de $p$, é necessário e suficiente que torne invertíveis as flechas pertencentes a $S$, e então a fatoração é única. A unicidade de tal par salvo isomorfismo é clara. Remetemos o leitor à exposição I do Seminário Homotópico de Estrasburgo para a demonstração da existência, bem como para a maioria dos resultados enunciados sem demonstração no que segue do número. Veja-se também P. Gabriel e M. Zisman, \emph{Calculus of fractions and homotopy theory}, Erg. der Math., 35, Springer-Verlag, 1967.

As propriedades do subconjunto multiplicativo $S$ fornecem uma descrição simples da categoria $\cE(S^{-1})$.

\begin{statement}{Definição 2.3.1}
Diz-se que um subconjunto $S$ de $\Fl(\cE)$ permite um cálculo de frações pela direita se as condições seguintes forem satisfeitas:
\begin{enumerate}[label=(\alph*)]
\item As flechas identidade de $\cE$ pertencem a $S$.
\item Se $u:X\to Y$ e $v:Y\to Z$ pertencem a $S$, a composição $v\circ u$ pertence a $S$.
\item Qualquer diagrama do tipo $(D_1)$ que aparece a seguir, em que $t$ pertence a $S$, pode ser completado em um diagrama do tipo $(D_2)$, em que $s$ pertence a $S$:
\[
\begin{tikzcd}[column sep=large,row sep=large]
(D_1)& & Y' \arrow[d,"t"]\\
& X \arrow[r,"u"'] & Y
\end{tikzcd}
\qquad
\begin{tikzcd}[column sep=large,row sep=large]
(D_2)& X' \arrow[r,"u'"] \arrow[d,"s"'] & Y' \arrow[d,"t"]\\
& X \arrow[r,"u"'] & Y .
\end{tikzcd}
\]
\item Para todo diagrama comutativo
\[
\begin{tikzcd}[column sep=large]
X \arrow[r,shift left,"f"] \arrow[r,shift right,"g"'] & Y \arrow[r,"s"] & Z,
\end{tikzcd}
\]
onde $s$ pertence a $S$, existe um diagrama comutativo
\[
\begin{tikzcd}[column sep=large]
T \arrow[r,"t"'] & X \arrow[r,shift left,"f"] \arrow[r,shift right,"g"'] & Y,
\end{tikzcd}
\]
onde $t$ pertence a $S$.
\end{enumerate}
\end{statement}

De maneira dual, define-se a noção seguinte.

\begin{statement}{Definição 2.3.2}
Diz-se que um subconjunto $S$ de $\Fl(\cE)$ permite um cálculo de frações pela esquerda se as condições seguintes forem satisfeitas:
\begin{enumerate}[label=(\alph*)]
\item[(a')] As flechas identidade de $\cE$ pertencem a $S$.
\item[(b')] A composição de duas flechas pertencentes a $S$ pertence a $S$.
\item[(c')] Qualquer diagrama do tipo $(D_1)'$ que aparece a seguir, em que $s$ pertence a $S$, pode ser completado em um diagrama do tipo $(D_2)'$, em que $t$ pertence a $S$:
\[
\begin{tikzcd}[column sep=large,row sep=large]
(D_1)'& X' \arrow[r,"u'"] \arrow[d,"s"'] & Y'\\
& X
\end{tikzcd}
\qquad
\begin{tikzcd}[column sep=large,row sep=large]
(D_2)'& X' \arrow[r,"u'"] \arrow[d,"s"'] & Y' \arrow[d,"t"]\\
& X \arrow[r,"u"'] & Y .
\end{tikzcd}
\]
\item[(d)'] Para todo diagrama comutativo $\Delta$ da forma que se apresenta a seguir, no qual $s$ pertence a $S$, existe um morfismo $t$ com fonte $Y$ e pertencente a $S$ tal que $t\circ f=t\circ g$:
\[
(\Delta)\quad
\begin{tikzcd}[column sep=large]
X' \arrow[r] & X \arrow[r,shift left,"f"] \arrow[r,shift right,"g"'] & Y .
\end{tikzcd}
\]
\end{enumerate}
\end{statement}
